\section{监督微调（SFT）：数据与方法}

监督微调是后训练流程的第一步，也是最直观的一步：收集高质量的``指令-回答''对，然后在预训练模型上做梯度下降。但看似简单的过程中，数据的选择和质量控制却是决定成败的关键。

\subsection{指令微调数据的三种范式}

课上详细对比了三个代表性数据集，它们分别代表三种不同的数据构建范式。

\begin{figure}[H]
\centering
\includegraphics[width=0.95\textwidth]{slides-images/slide-008.jpg}
\caption{三种指令微调数据集的对比：FLAN（NLP 任务聚合）、Stanford Alpaca（LLM 生成）和 Open Assistant（人类众包）\protect\footnotemark}
\end{figure}
\footnotetext{Slides 第9页。}

\subsubsection{FLAN：NLP 任务聚合}

FLAN 由 Google 团队构建，其策略是将大量已有的 NLP 数据集（如 Natural Instructions V2、T0-SF、对抗性 QA、主题分类等）聚合成一个大的元数据集。每个数据集都被转换为``指令+输入+输出''的格式。

\begin{knowledgebox}{FLAN 的优劣势}
\textbf{优势}：数据量大、覆盖任务类型广泛、构建成本低（复用已有数据集）。\\
\textbf{劣势}：任务格式不自然（如多选题、Enron 邮件主题生成等）；输出通常很短（一个词或短语）；与真实的聊天交互差异较大。
\end{knowledgebox}

\begin{figure}[H]
\centering
\includegraphics[width=0.95\textwidth]{slides-images/slide-009.jpg}
\caption{FLAN 数据集随机样本：包含文章摘要、多选分类、邮件主题生成、数据库到文本等多种 NLP 任务\protect\footnotemark}
\end{figure}
\footnotetext{Slides 第10页。}

\subsubsection{Stanford Alpaca：LLM 生成数据}

Alpaca 是早期使用语言模型生成指令微调数据的代表性工作。其流程是：用一组人工编写的种子指令作为上下文示例，让语言模型生成新的指令，再用 InstructGPT 生成对应的回复。

\begin{figure}[H]
\centering
\includegraphics[width=0.95\textwidth]{slides-images/slide-010.jpg}
\caption{Alpaca 数据集随机样本：指令更接近日常聊天风格，回复为长文本自然语言\protect\footnotemark}
\end{figure}
\footnotetext{Slides 第11页。}

与 FLAN 相比，Alpaca 的指令更接近用户可能输入到聊天机器人的内容，回复也更偏向长文本自然语言。但其输入多样性有限（指令往往较短且简单），且完全依赖 LLM 的生成质量。

\subsubsection{Open Assistant：人类众包}

Open Assistant 是 ChatGPT 发布后，由一群在线爱好者自发组织的众包项目。志愿者们编写了大量指令微调数据，其特点是：

\begin{itemize}
    \item 查询更复杂、更贴近真实需求
    \item 回复非常详细，甚至包含引用文献
    \item 数据质量参差不齐，但优质样本质量很高
\end{itemize}

\begin{figure}[H]
\centering
\includegraphics[width=0.95\textwidth]{slides-images/slide-011.jpg}
\caption{Open Assistant 数据集随机样本：展示了更复杂的查询和包含引用的详细回复\protect\footnotemark}
\end{figure}
\footnotetext{Slides 第12页。}

\subsection{数据质量的关键维度}

课上通过课堂互动实验揭示了指令微调数据面临的核心挑战。

\subsubsection{长度偏差}

\begin{figure}[H]
\centering
\includegraphics[width=0.95\textwidth]{slides-images/slide-015.jpg}
\caption{不同指令微调数据集的输入/输出长度分布差异极大\protect\footnotemark}
\end{figure}
\footnotetext{Slides 第16页。来自 Wang et al. 2023 的综述。}

长度是后训练数据中一个``房间里的大象''：

\begin{warningbox}{长度偏差的危害}
\begin{itemize}
    \item 人类评估者对列表格式和较长输出有强烈偏好（约 60--70\% 的偏好率）
    \item AI 评审（如 GPT-4 作为 judge）同样存在此偏差
    \item 这意味着仅通过增加输出长度就可以``欺骗''评估系统，而非真正提升内容质量
\end{itemize}
一个好的后训练流程应该优化实质性内容（如减少幻觉、提升准确率），而非风格因素。
\end{warningbox}

有趣的是，这些长度相关因素对基准测试（如 MMLU）的影响较小——大多数指令微调数据集都能在基准测试上带来相对于基座模型的提升。因此，在评估后训练效果时，应当综合使用聊天式评估（如 Chatbot Arena、AlpacaEval）和标准基准测试，以避免长度偏差的干扰。

\subsubsection{知识注入与幻觉的权衡}

这是本课中最深刻的洞见之一。考虑 Open Assistant 中的一个示例：用户询问``单一消费者经济学''（monopsony），回复中包含了具体的学术引用。

\begin{figure}[H]
\centering
\includegraphics[width=0.95\textwidth]{slides-images/slide-017.jpg}
\caption{SFT 数据中的知识注入问题：当回复包含模型可能不知道的引用时，微调同时教授两件事\protect\footnotemark}
\end{figure}
\footnotetext{Slides 第18页。}

当我们用这样的数据微调模型时，模型同时学到两件事：

\begin{importantbox}{SFT 中的两种竞争机制}
\begin{enumerate}
    \item \textbf{知识关联}（正面）：将``单一消费者经济学''与特定引用关联起来
    \item \textbf{格式模仿}（潜在负面）：学到``遇到复杂概念时，应在末尾添加参考文献''的模式
\end{enumerate}
如果模型参数中本就不包含该引用的知识，那么第二种机制会占主导——模型学会的不是检索正确引用，而是\textbf{编造一个看似合理的引用}。
\end{importantbox}

John Schulman（OpenAI）曾在 Berkeley 的报告中强调：如果模型不具备回答某问题的知识，强迫它回答只会教它幻觉。正确的做法是让模型学会说``我不知道''。

\begin{warningbox}{反直觉的数据质量准则}
一个``完全正确、信息丰富''的微调数据集，可能反而对模型有害——因为当数据的知识深度超出模型预训练时学到的内容时，模型会学到``编造事实以匹配回复格式''的捷径行为。这一问题在以下两种情况下尤其突出：
\begin{itemize}
    \item 蒸馏数据中教师模型远强于学生模型
    \item 人类标注者的知识远超模型的能力
\end{itemize}
\end{warningbox}

\subsection{安全微调}

\begin{figure}[H]
\centering
\includegraphics[width=0.95\textwidth]{slides-images/slide-022.jpg}
\caption{安全微调的核心权衡：在拒绝有害请求和避免过度拒绝之间找到平衡\protect\footnotemark}
\end{figure}
\footnotetext{Slides 第23页。}

安全微调面临一个核心权衡：模型需要拒绝有害请求（如``写一条仇恨言论的推文''），但不能过度拒绝合理请求（如``如何 kill 一个 Python 进程''）。

关键发现包括：
\begin{itemize}
    \item 少量安全微调数据（甚至 500 个样本）就能显著提升模型的安全性
    \item 这与指令微调的发现一致：强预训练模型 + 少量高质量数据 = 显著的行为改变
    \item 但纯粹通过 SFT 来处理安全问题的细微之处非常困难
\end{itemize}

\subsection{SFT 的方法论}

\subsubsection{基本方法：梯度下降}

SFT 的基本操作非常直接：将指令和回复拼接为序列，对回复部分计算交叉熵损失，做梯度下降。在学术场景中，这通常就够了。

\subsubsection{规模化：中间训练（Mid-training）}

\begin{figure}[H]
\centering
\includegraphics[width=0.95\textwidth]{slides-images/slide-025.jpg}
\caption{MiniCPM 的两阶段训练流程：第一阶段为纯预训练，第二阶段（衰减阶段）混入指令微调数据\protect\footnotemark}
\end{figure}
\footnotetext{Slides 第26页。来自 MiniCPM 论文。}

在前沿实验室中，现代指令微调流程正在与预训练流程融合。具体做法是：

\begin{importantbox}{中间训练范式}
\begin{enumerate}
    \item \textbf{第一阶段}：纯预训练（使用 Common Crawl、代码数据等）
    \item \textbf{第二阶段}（衰减/退火阶段）：在学习率下降的同时，逐步混入指令微调数据（如 UltraChat、Stack Exchange QA、代码 SFT 数据等）
    \item \textbf{第三阶段}（可选）：在第二阶段完成后，进行一个短的纯指令微调轮次
\end{enumerate}
\end{importantbox}

以 MiniCPM 为例，其第一阶段使用纯预训练数据（Common Crawl、代码预训练、Pile、Dolma 等），第二阶段在学习率衰减时混入代码 SFT、中文书籍、UltraChat、Stack Exchange QA 等指令微调相关数据。

这种范式的优势包括：

\begin{itemize}
    \item 避免灾难性遗忘（因为预训练数据仍在混合中）
    \item 指令微调数据能更深入地整合到模型中
    \item 可以利用学习率衰减带来的大幅损失下降来``退火''模型进入正确的模式
\end{itemize}

\begin{warningbox}{``基座模型''概念的模糊化}
在这种范式下，所谓的``基座模型''实际上已经通过中间训练阶段隐式地接受了指令微调。当我们看到 Qwen 等公司发布的``base model''时，它很可能已经经历了包含指令数据的衰减阶段。因此，``基座模型''这个术语正变得越来越具有歧义性。
\end{warningbox}

\subsection{本章小结}

\begin{importantbox}{SFT 部分的三个核心要点}
\begin{enumerate}
    \item \textbf{指令微调惊人地强大}：即使使用标准的开源数据集（如 Open Hermes、Open Assistant），配合合理的超参数，就能得到一个行为类似 LLaMA 或 ChatGPT 的模型
    \item \textbf{``高质量数据''的定义非常复杂}：完全正确的数据不一定是好数据；需要考虑模型的知识边界、幻觉风险等多个维度
    \item \textbf{少量数据即可产生巨大影响}：500 个安全样本就能显著改变模型的安全行为
\end{enumerate}
\end{importantbox}


